\section{练习}

\begin{question}\begin{drill}[横档阻值改变的双向梯形]\label{c4:drill:1}
双向无限双导轨每段导轨电阻 $r$，横档均为 $R$。求同一横档两端电阻，并讨论 $R\to0$ 和 $R\to\infty$。
\end{drill}
\typeout{ECFIG:fig:c4:drill:1:s}\fig{fig:c4:drill:1:s}{题设：横档阻值改变的双向梯形}{\begin{tikzpicture}[scale=.88]
\draw(0.0,0)to[R,l={$R$}](0.0,2);\draw(0.0,2)to[R,l={$r$}](1.7,2);\draw(0.0,0)to[R,l_={$r$}](1.7,0);\draw(1.7,0)to[R,l={$R$}](1.7,2);\draw(1.7,2)to[R,l={$r$}](3.4,2);\draw(1.7,0)to[R,l_={$r$}](3.4,0);\draw(3.4,0)to[R,l={$R$}](3.4,2);\draw(3.4,2)to[R,l={$r$}](5.1,2);\draw(3.4,0)to[R,l_={$r$}](5.1,0);\draw(5.1,0)to[R,l={$R$}](5.1,2);\draw[dashed](5.1,0)--(6.1,0)(5.1,2)--(6.1,2);\node[left]at(0,2){$A$};\node[left]at(0,0){$B$};\draw[dashed](-1,0)--(0,0)(-1,2)--(0,2);
\end{tikzpicture}}
\end{question}
\answerlink{c4:drill:1}
\begin{question}\begin{drill}[无限方格去掉直接边]\label{c4:drill:2}
无限方格每边为 $r$，相邻端点 $A,B$ 间的直接电阻断开，其余网络不变。求新的端口电阻。
\end{drill}
\typeout{ECFIG:fig:c4:drill:2:s}\fig{fig:c4:drill:2:s}{题设：无限方格去掉直接边}{\begin{tikzpicture}[scale=.88]
\draw[world](0,-2)--(0,2);\draw[world](0,-2)--(4,-2);\draw[world](1,-2)--(1,2);\draw[world](0,-1)--(4,-1);\draw[world](2,-2)--(2,2);\draw[world](0,0)--(4,0);\draw[world](3,-2)--(3,2);\draw[world](0,1)--(4,1);\draw[world](4,-2)--(4,2);\draw[world](0,2)--(4,2);\fill(1,0)circle(2pt)node[below left]{$A$};\fill(2,0)circle(2pt)node[below right]{$B$};\node at(2,2.7){每边 $r$，向四周无限延伸};\draw[white,line width=3pt](1.18,0)--(1.82,0);\draw[guide](1.18,0)--(1.82,0);
\end{tikzpicture}}
\end{question}
\answerlink{c4:drill:2}
\begin{question}\begin{drill}[无限蜂窝的一条边改变]\label{c4:drill:3}
无限蜂窝每边原为 $r$，仅相邻端口 $A,B$ 的直接边改为 $R$。求等效电阻。
\end{drill}
\typeout{ECFIG:fig:c4:drill:3:s}\fig{fig:c4:drill:3:s}{题设：无限蜂窝的一条边改变}{\begin{tikzpicture}[scale=.88]
\draw[world](-0.5,-0.8660254037844386)--(-0.9999999999999999,0.0);\draw[world](-0.9999999999999999,0.0)--(-1.9999999999999998,1.1102230246251565e-16);\draw[world](-1.9999999999999998,1.1102230246251565e-16)--(-2.5,-0.8660254037844385);\draw[world](-2.5,-0.8660254037844385)--(-2.0000000000000004,-1.732050807568877);\draw[world](-2.0000000000000004,-1.732050807568877)--(-0.9999999999999999,-1.7320508075688772);\draw[world](-0.9999999999999999,-1.7320508075688772)--(-0.5,-0.8660254037844386);\draw[world](1.0,-1.7320508075688772)--(0.5000000000000001,-0.8660254037844386);\draw[world](0.5000000000000001,-0.8660254037844386)--(-0.49999999999999983,-0.8660254037844385);\draw[world](-1.0,-1.732050807568877)--(-0.5000000000000004,-2.5980762113533156);\draw[world](-0.5000000000000004,-2.5980762113533156)--(0.5000000000000001,-2.598076211353316);\draw[world](0.5000000000000001,-2.598076211353316)--(1.0,-1.7320508075688772);\draw[world](2.5,-0.8660254037844386)--(2.0,0.0);\draw[world](2.0,0.0)--(1.0000000000000002,1.1102230246251565e-16);\draw[world](1.0000000000000002,1.1102230246251565e-16)--(0.5,-0.8660254037844385);\draw[world](0.9999999999999996,-1.732050807568877)--(2.0,-1.7320508075688772);\draw[world](2.0,-1.7320508075688772)--(2.5,-0.8660254037844386);\draw[world](-0.5,0.8660254037844386)--(-0.9999999999999999,1.7320508075688772);\draw[world](-0.9999999999999999,1.7320508075688772)--(-1.9999999999999998,1.7320508075688772);\draw[world](-1.9999999999999998,1.7320508075688772)--(-2.5,0.8660254037844387);\draw[world](-2.5,0.8660254037844387)--(-2.0000000000000004,2.220446049250313e-16);\draw[world](-0.9999999999999999,0.0)--(-0.5,0.8660254037844386);\draw[world](1.0,0.0)--(0.5000000000000001,0.8660254037844386);\draw[world](0.5000000000000001,0.8660254037844386)--(-0.49999999999999983,0.8660254037844387);\draw[world](2.5,0.8660254037844386)--(2.0,1.7320508075688772);\draw[world](2.0,1.7320508075688772)--(1.0000000000000002,1.7320508075688772);\draw[world](1.0000000000000002,1.7320508075688772)--(0.5,0.8660254037844387);\draw[world](2.0,0.0)--(2.5,0.8660254037844386);\draw[world](-0.5,2.598076211353316)--(-0.9999999999999999,3.4641016151377544);\draw[world](-0.9999999999999999,3.4641016151377544)--(-1.9999999999999998,3.464101615137755);\draw[world](-1.9999999999999998,3.464101615137755)--(-2.5,2.598076211353316);\draw[world](-2.5,2.598076211353316)--(-2.0000000000000004,1.7320508075688776);\draw[world](-0.9999999999999999,1.7320508075688774)--(-0.5,2.598076211353316);\draw[world](1.0,1.7320508075688772)--(0.5000000000000001,2.598076211353316);\draw[world](0.5000000000000001,2.598076211353316)--(-0.49999999999999983,2.598076211353316);\draw[world](2.5,2.598076211353316)--(2.0,3.4641016151377544);\draw[world](2.0,3.4641016151377544)--(1.0000000000000002,3.464101615137755);\draw[world](1.0000000000000002,3.464101615137755)--(0.5,2.598076211353316);\draw[world](2.0,1.7320508075688774)--(2.5,2.598076211353316);\fill(-.5,-.866)circle(1.7pt)node[below left]{$A$};\fill(.5,-.866)circle(1.7pt)node[below right]{$B$};\fill(1,0)circle(1.7pt)node[right]{$C$};\draw[white,line width=3pt](-.5,-.866)--(.5,-.866);\draw(-.5,-.866)to[R,l={$R$}](.5,-.866);
\end{tikzpicture}}
\end{question}
\answerlink{c4:drill:3}
\begin{question}\begin{drill}[有限梯形末端边界的作用]\label{c4:drill:4}
每节先串联 $r$、再接横档 $r$ 的单导轨梯形，只保留两节。分别令最末横档之后开路及接理想短路，求输入电阻；讨论是否可以直接使用无限固定点。
\end{drill}
\typeout{ECFIG:fig:c4:drill:4:s}\fig{fig:c4:drill:4:s}{题设：有限梯形末端边界的作用}{\begin{tikzpicture}[scale=.88]
\draw(0.0,2)to[R,l={$r$}](1.8,2);\draw(1.8,2)to[R,l={$r$}](1.8,0);\draw(1.8,2)to[R,l={$r$}](3.6,2);\draw(3.6,2)to[R,l={$r$}](3.6,0);\draw(0,0)--(3.6,0);\node[left]at(0,2){$A$};\node[left]at(0,0){$B$};
\end{tikzpicture}}
\end{question}
\answerlink{c4:drill:4}
\begin{question}\begin{drill}[单导轨的电阻与端口间距]\label{c4:drill:5}
双向单导轨网络每电阻 $r$、下导线理想。两上轨节点相距两段。求电阻，并说明相距越来越远时电阻为何趋有限值。
\end{drill}
\typeout{ECFIG:fig:c4:drill:5:s}\fig{fig:c4:drill:5:s}{题设：单导轨的电阻与端口间距}{\begin{tikzpicture}[scale=.88]
\draw(0.0,0)to[R,l={$r$}](0.0,2);\draw(0.0,2)to[R,l={$r$}](1.7,2);\draw(0.0,0)--(1.7,0);\draw(1.7,0)to[R,l={$r$}](1.7,2);\draw(1.7,2)to[R,l={$r$}](3.4,2);\draw(1.7,0)--(3.4,0);\draw(3.4,0)to[R,l={$r$}](3.4,2);\draw(3.4,2)to[R,l={$r$}](5.1,2);\draw(3.4,0)--(5.1,0);\draw(5.1,0)to[R,l={$r$}](5.1,2);\draw[dashed](5.1,0)--(6.1,0)(5.1,2)--(6.1,2);\node[left]at(0,2){$A_0$};\draw[dashed](-1,0)--(0,0)(-1,2)--(0,2);
\fill(3.4,2)circle(1.8pt)node[above]{$A_2$};\end{tikzpicture}}
\end{question}
\answerlink{c4:drill:5}
\begin{question}\begin{drill}[双导轨相邻异轨端口]\label{c4:drill:6}
双向无限双导轨每边为 $r$。$A$ 在零列上轨，$D$ 在第一列下轨。求电阻，并推广到相隔 $n$ 列的异轨端口。
\end{drill}
\typeout{ECFIG:fig:c4:drill:6:s}\fig{fig:c4:drill:6:s}{题设：双导轨相邻异轨端口}{\begin{tikzpicture}[scale=.88]
\draw(0.0,0)to[R,l={$r$}](0.0,2);\draw(0.0,2)to[R,l={$r$}](1.7,2);\draw(0.0,0)to[R,l_={$r$}](1.7,0);\draw(1.7,0)to[R,l={$r$}](1.7,2);\draw(1.7,2)to[R,l={$r$}](3.4,2);\draw(1.7,0)to[R,l_={$r$}](3.4,0);\draw(3.4,0)to[R,l={$r$}](3.4,2);\draw(3.4,2)to[R,l={$r$}](5.1,2);\draw(3.4,0)to[R,l_={$r$}](5.1,0);\draw(5.1,0)to[R,l={$r$}](5.1,2);\draw[dashed](5.1,0)--(6.1,0)(5.1,2)--(6.1,2);\node[left]at(0,2){$A$};\node[left]at(0,0){$B$};\draw[dashed](-1,0)--(0,0)(-1,2)--(0,2);\node[below]at(1.7,0){$D$};
\end{tikzpicture}}
\end{question}
\answerlink{c4:drill:6}
